\chapter{有机合成像搭积木，也像规划路线}

\begin{kaichang}
去药箱里翻一翻，很可能能找到一板阿司匹林：白色的药片，说明书上印着一个拗口的名字——乙酰水杨酸。全世界每年大约要生产四万吨这种东西，堆起来能装满十几节火车皮。

再想象另一样东西：一杯柳树皮煮的水，棕黄色，苦得难以下咽。两千多年前的古希腊人就喝它退烧止痛，中国的古人也用柳枝“去热”。

现在请你在纸上写下一个答案：如果这两样东西都能退烧，你选哪一个？顺便写下理由。我猜很多人会选柳树皮水——“天然的，总比化工厂里造出来的安全吧？”

这一章结束时，我们回来对答案。你会发现，这个问题根本没有标准答案里那一栏——因为“天然”和“合成”这两个标签，从一开始就问错了方向。真正的问题藏在分子层面：它是哪个分子？怎么造出来的？造它的时候产生了多少别的东西？
\end{kaichang}

\section{一颗药片的两种“出身”}

先看能看见、能测到的东西。

阿司匹林片碾碎了是白色粉末，微微发酸，几乎不溶于冷水。柳树皮煮的水是苦的——苦味来自其中一类叫水杨苷的物质。十九世纪的化学家把水杨苷从柳树皮里分离出来，发现它进入人体后会转变成另一个分子：水杨酸。水杨酸确实能退烧止痛，但它酸性强、对胃的刺激大，直接吃很难受。

于是化学家做了一件事：给水杨酸分子动一个小“手术”——在它的羟基（\chem{-OH}）上加一顶乙酰基“小帽子”。改造后的分子就是乙酰水杨酸，也就是阿司匹林。改造的目的很朴素：让分子更便于服用、对胃肠的耐受好一些。注意，\textbf{分子骨架几乎没动}，药效的根基还是那个水杨酸骨架。

这里出现一个贯穿本章的问题：这个“手术”是怎么做的？化学家面对一个目标分子，怎样知道该从什么原料出发、先做什么、后做什么？

在回答之前，先记住一个事实：用第 35 章学过的化学侦探术（熔点、色谱、光谱）去检验，工厂造出的乙酰水杨酸和理论上能从天然原料出发制得的乙酰水杨酸，\textbf{是同一个分子}。分子不记日记，它身上没有任何标记能说明自己“出身”于柳树皮还是反应釜。这一点后面还要回来细谈。

\section{先学会拆，再学会搭}

假设你是化学家，老板给你一个目标分子：乙酰水杨酸。你怎么办？

直接的想法是“试”：把各种原料混在一起加热，看看能不能撞上产物。十九世纪早期确实有人这么干，成功率嘛——大致相当于把一堆零件抛向空中，指望它们落地时自己组装成一块手表。

化学家后来学会了一种反着来的思路：\textbf{先拆后搭}。盯着目标分子看，问：如果把这根键断开，会得到哪两个更简单的碎片？这两个碎片各自又能不能再拆？一路拆下去，直到拆成货架上买得到的便宜原料为止。拆的过程是倒着走的，搭的过程才是顺着走的。这套“倒着想”的系统方法叫\textbf{逆合成分析}，是美国化学家科里在二十世纪六十年代整理成形的——但他反复强调，这不是某个天才的灵感，而是一套谁都能学的思维工具：像修表师傅面对一块表，第一件事不是拧螺丝，而是看懂图纸，想清楚它是怎么装起来的。

\begin{center}
\begin{tikzpicture}[scale=1.0]
\node[draw,rectangle,minimum width=2.6cm,minimum height=0.9cm] (a) at (0,0) {便宜原料};
\node[draw,rectangle,minimum width=2.6cm,minimum height=0.9cm] (b) at (4.2,0) {中间体};
\node[draw,rectangle,minimum width=2.6cm,minimum height=0.9cm] (c) at (8.4,0) {目标分子};
\draw[->,thick] (a) -- (b);
\draw[->,thick] (b) -- (c);
\draw[->,dashed,thick] (c.north) .. controls (6.3,1.1) .. (b.north);
\draw[->,dashed,thick] (b.north) .. controls (2.1,1.1) .. (a.north);
\node at (4.2,-0.9) {\small 实线：实际合成方向};
\node at (4.2,1.6) {\small 虚线：倒推分析方向（逆合成）};
\end{tikzpicture}
\end{center}

这幅图是人为的表示法：方块和箭头只是帮助我们整理思路，不代表分子的真实模样。

拿阿司匹林练练手。盯着乙酰水杨酸看，它身上有一个酯的结构（第 40 章讲过，酯是羧酸和醇“拉手”形成的，拉手处叫酯键）。逆合成分析的第一刀就切在这里：把酯键断开，得到两个碎片——水杨酸和一个提供乙酰基的小分子。水杨酸买得到，乙酰基可以用乙酸酐（你可以把它粗略理解成“两个乙酸分子挤掉一分子水后的成品”）提供。好了，拆到头了：两种原料都便宜易得。整个阿司匹林的“图纸”，一层就拆完了。

当然，大多数药物分子比阿司匹林复杂得多，要拆十几层甚至几十层，画出来的不是一条线，而是一棵倒着长的树。思路完全一样：每一刀都要切在“断得开、也接得回去”的键上。

\section{三个“在哪里”的问题}

图纸有了，麻烦才刚开始。真实分子身上往往不止一个能反应的位置，就像一块电路板上不止一个焊点。每动一次手，化学家要同时回答三个问题：

\begin{itemize}[itemsep=2pt]
\item \textbf{在哪个位置反应？}分子上有好几个官能团（第 37 章），试剂凭什么只动其中一个、不去碰别的？这叫化学选择性。比如水杨酸身上既有羟基又有羧基，加乙酰基时只想动羟基——选错试剂，羧基也会跟着起哄。
\item \textbf{反应到哪一步停下来？}第 40 章讲过，醇可以被氧化成醛，醛还能继续被氧化成羧酸。想要醛，就得让反应在半路上刹车——选温和的氧化剂，或者让醛一生成就赶紧离开反应体系。
\item \textbf{生成哪个立体结构？}第 42 章讲过手性：同一个分子有左手和右手两种镜像版本，药效可能天差地别。反应造出新的手性中心时，往往左右手一起生成，怎么让想要的版本占上风？
\end{itemize}

这三个“在哪里”回答不好，图纸再漂亮也白搭。那回答不好的时候怎么办？化学家发明了一个笨但有效的办法：\textbf{保护基}。

想象你要给墙壁刷漆，但墙上有一个开关面板不想沾漆。最直接的办法是什么？拿美纹纸把开关贴住，刷完漆再撕掉。保护基就是化学世界的美纹纸：想动 A 基团、又怕 B 基团跟着反应，就先给 B 戴一顶“保护帽”（上一个保护基），等 A 那边的活干完了，再把帽子摘掉。

代价显而易见：多出来“戴帽子”和“摘帽子”两步反应，多花原料，多产生废物，总产率还要再打折。保护基是“必要的浪费”——路线设计者心里都清楚，能不用就不用。怎样才算“能不用”？那就得找天生就挑位置的试剂。这个念头先记住，章末它会变成一颗重要的种子。

\section{每一步都要交“路费”}

现在谈钱——不，谈产率。没有一步化学反应是百分之百转化的：总有一些原料没反应，总有一些产物在分离提纯时损失掉。每步反应的收获比例叫产率。

假设一条合成路线有十步，每步产率都是相当体面的 $90\%$。总产率是多少？直觉会说“大概还是九成吧”。算一下：

\[0.9^{10}\approx 0.35\]

只剩百分之三十五。投进去一百克原料，最终只拿到三十五克产品，其余六十五克都变成了各步的“路费”。这就是多步合成的残酷数学：\textbf{产率是连乘的}，每一步的损耗都在放大下一步的损耗。二十步以后，$0.9^{20}\approx 0.12$，八成以上的原料都打了水漂。

\begin{qiaoliang}
产率之外，还有一笔更根本的账：\textbf{原子经济性}。它问的不是“我投进去的原料回收了多少产品”，而是“反应物的原子里，有多少最终留在了我想要的产物分子里”。计算方法是：目标产物的摩尔质量，除以所有反应物摩尔质量之和。

拿阿司匹林这一步来算。水杨酸（\chem{C_7H_6O_3}，摩尔质量 \ud{138}{g\cdot mol^{-1}}）和乙酸酐（\chem{C_4H_6O_3}，\ud{102}{g\cdot mol^{-1}}）反应，生成乙酰水杨酸（\chem{C_9H_8O_4}，\ud{180}{g\cdot mol^{-1}}）和一分子乙酸（\ud{60}{g\cdot mol^{-1}}）：

\[\text{原子经济性}=\frac{180}{138+102}=\frac{180}{240}=75\%\]

也就是说，即使这步反应做到完美（产率 $100\%$），投进去的原子里仍有四分之一注定要变成副产物乙酸。原子经济性是路线的“先天体质”，产率是“后天发挥”——体质差的路线，发挥再好也有废物上限。绿色化学的核心思想就藏在这笔账里：不要等废物产生了再治理，而要在画图纸的阶段就选择原子浪费少的路线。比如加成反应（第 38 章，两个分子合成一个、不产生副产物）的原子经济性天然就是 $100\%$，是路线设计师的最爱。
\end{qiaoliang}

\begin{shiyan}
\textbf{掷骰子模拟一条十步合成路线}（纯纸笔模拟，不使用任何化学品）。

材料：一颗骰子、纸和笔。

步骤：假设你在做一条十步的合成路线，每步的产率是六分之五（约 $83\%$）——用骰子模拟：掷出 1 到 5 算这步成功，掷出 6 算这步失败、整锅报废。规则：从第 1 步开始掷，连续成功走完 10 步才算拿到最终产品；中途掷出 6 就从头再来。

观察点：玩 20 局，记录有几局成功走到终点。理论上成功的概率是 $(5/6)^{10}\approx 16\%$，也就是 20 局里大约只能成功 3 局。你的结果接近吗？再把规则改成每步成功率 $35/36$（掷出 6 允许重掷一次，再出 6 才算失败，成功率约 $97\%$），感受“每步提高一点点，全程相差巨大”。

安全提示：这个实验只用纸笔骰子，没有任何危险。但请记住：真实的化学合成涉及有毒溶剂、强腐蚀性试剂和高温玻璃仪器，\textbf{绝对不要在家模仿任何真实的合成操作}，想看真实的反应釜，去搜正规的科普视频或等老师的课堂演示。
\end{shiyan}

有了产率和原子经济性这两把尺子，就能看懂工业界的真实选择了。止痛药水杨酸的“后辈”布洛芬是个经典案例：老工艺要六步，原子经济性只有约 $40\%$；二十世纪九十年代的新工艺砍到三步，原子经济性升到约 $77\%$（如果把副产物回收利用，可接近 $99\%$），废物量砍掉一大半，还因此拿了绿色化学领域的大奖。注意，新工艺赢的秘诀不是更猛的反应条件，而是更聪明的路线设计——这正是本章的主题。

\section{让反应式登场}

前面故意忍着没写反应式。现在符号可以登场了——阿司匹林的那一步：

\[\chem{C_7H_6O_3} + \chem{C_4H_6O_3} \rightarrow \chem{C_9H_8O_4} + \chem{C_2H_4O_2}\]

水杨酸加乙酸酐，生成乙酰水杨酸，附带一分子乙酸。数一数两边的原子：碳、氢、氧的数目完全相等——原子守恒（第 24 章）在这里默默站岗。这行式子背后是前面所有的讨论：逆合成决定了“为什么选这两种原料”，选择性决定了“为什么只动羟基”，原子经济性告诉我们那分子乙酸是躲不掉的副产品。

再强调一次读法：箭头不是“凭魔法变成”，而是一大批分子按第 38 章讲的机理（电子对的移动）一步步重新排列的统计结果；反应能朝这个方向走，受第 27 章能量规则的约束；走多快，受第 28 章速率规则的约束。图纸、能量、速率，三者缺一不可。

\section{从霉菌到抗生素：一条路线的百年接力}

阿司匹林的路线只有一步，太简单了，不足以展示“规划路线”真正的威力。讲一个更硬的故事。

\begin{zhengju}
1928 年，伦敦的细菌学家弗莱明度假回来，发现一个被霉菌污染的培养皿有些古怪：青绿色的霉菌斑点周围，葡萄球菌都不长了。第一反应可能是“霉菌把营养抢光了吧”——这是一个合理的替代解释。弗莱明做了一个关键实验：把霉菌培养过的液体过滤，\textbf{滤液里已经没有活的霉菌了}，可它照样能杀死细菌。这排除了“抢营养”，指向一个新结论：霉菌向环境里释放了一种能杀菌的物质。他把它叫作青霉素。

但弗莱明卡在了下一步：这种物质极不稳定，他提纯不出来，浓度低得没法用于治病。他甚至一度认为青霉素最多只能当外用的局部消毒剂。发现不等于能用——中间的鸿沟，要靠“制造路线”来填。

十年后，牛津大学的弗洛里和钱恩团队接过了接力棒。他们用冷冻干燥等技术拿到了较纯的青霉素，1940 年做了那个著名的动物实验：8 只注射了致死量病菌的小鼠，4 只用青霉素治疗，活了下来；4 只不治疗，全部死亡。证据够硬，可新的问题来了：治疗一个人需要的青霉素，要从几百升霉菌培养液里提取。战时需求逼人，科学家跑到发霉的甜瓜上筛到了产量更高的霉菌菌株，又改用大罐子深层发酵——\textbf{让霉菌替我们做合成的苦力}，产量才翻了成百上千倍。

还剩最后一个谜：青霉素分子到底长什么样？化学家吵了好几年，提出过不同的结构。1945 年，霍奇金用 X 射线晶体学（第 35 章的侦探术之一）拍板定案：青霉素的核心是一个紧张的四元环（$\beta$-内酰胺环）。结构一清楚，化学家立刻看懂了它的弱点——这个环太容易被酸和细菌产生的酶掰开。于是五十年代末，人们从发酵液里直接提取青霉素的“裸机底盘”（去掉侧链的母核），再用化学方法往底盘上接各种不同的侧链：抗酸的、抗菌谱更广的、不怕细菌酶的……一批“半合成青霉素”就这样从图纸上走了下来，至今仍在救人。

回看这条百年接力：发现现象、排除替代解释、提纯、定量证据、发酵量产、测定结构、改造分子——没有一环是某个天才的独舞，每一环都在回答本章的问题：分子从哪里来，怎样造得更多、更好、更合身。
\end{zhengju}

这个故事里藏着一个值得多看一眼的细节：青霉素直到今天也不靠“全合成”（从简单原料一步步拼出整个分子）来生产。不是化学家不会——1964 年就有人完成了它的全合成——而是全合成路线太长、太贵、产率太低，打不过“霉菌造底盘、化学家装修”的半合成模式。这直接引出下一个话题：这套图纸思维的边界在哪里。

\section{图纸不等于能造出来}

逆合成分析是一张神奇的地图，但地图不等于路。它有四条明显的边界：

\begin{itemize}[itemsep=2pt]
\item \textbf{纸面能断开的键，烧瓶里不一定接得上。}逆合成分析只管“在哪里断”，不管断开的两个碎片是否真能在现实条件下相遇成键。每一步都要同时过两关：能量关（第 27 章，方向是否被热力学允许）和速率关（第 28 章，活化能高不高、副反应快不快）。
\item \textbf{步数本身就是敌人。}前面算过连乘的账。复杂的天然分子动辄几十上百步，总产率常常不到百分之一。维生素 \chem{B_{12}} 的全合成就动员了两个团队、近百人、花了十一年——它的意义是证明“人类能做到”，而不是“应该这样生产”。真正生产的 \chem{B_{12}} 靠微生物发酵。
\item \textbf{路线图不告诉你哪条路最省。}逆合成能给出好几条候选路线，选哪条要算产率、原子经济性、原料价格、废物处理、安全性——这是工程和经济问题，超出了图纸本身。
\item \textbf{有些“选择性”几乎无解。}想让反应只发生在某个位置、只生成某一只手性的产物，化学试剂常常做不到，只能靠保护基硬扛，或者靠拆分扔掉一半产物。这是传统合成最深的痛处——记住它，章末的伏笔从这里长出来。
\end{itemize}

所以别被“能画出逆合成树”迷惑：图纸是思考的起点，可行性是试出来的。好的合成化学家一半时间在画图纸，另一半时间在烧瓶前被现实教育。

\begin{wuqu}
“天然的就是安全的，合成的就是有害的。”——这是本章最该拆掉的一枚标签。

反例从两边夹击。先看“天然等于安全”：蓖麻毒素、乌头碱、毒鹅膏蘑菇里的鹅膏毒素，全是百分之百的天然产物，毒性名列前茅；本章开头的水杨酸也是天然的，照样刺激胃。再看“合成等于有害”：合成的胰岛素、合成的抗生素每年救的人以百万计；而维生素 \chem{C} 泡腾片里的维生素 \chem{C} 是工厂合成的，柠檬里的维生素 \chem{C} 是植物合成的——两者是\textbf{同一个分子}，结构一模一样，性质一模一样，没有任何仪器能分出它们的“出身”。

毒性从来不由“出身”决定，而由\textbf{分子本身和剂量}决定（第 35 章说过：检出不等于有害，剂量才是关键）。用“天然/合成”给物质贴好坏标签，就像按衣服的产地判断一个人的人品——方便是方便，错也是真的错。
\end{wuqu}

\section{回到那板阿司匹林}

现在可以对开场的问题交卷了。柳树皮水和阿司匹林片，选哪个？

正确的回答是：这个问题一开始就问错了。有效成分是分子，不是出身。柳树皮水的麻烦不在于它“天然”，而在于\textbf{成分不明、剂量不可控}——这棵柳树和那棵柳树的水杨苷含量不同，煮的时间长短不同，你根本不知道一杯下去吃进了多少有效分子，还顺带喝进了一堆未经检验的杂质。阿司匹林片的优势也不在于它“合成”，而在于\textbf{成分明确、纯度有标准、每片的剂量精确}。制药工业的全部努力——正是本章讲的路线设计、纯化、质量控制——本质上就是为了让“每一颗药片里的分子种类和数量”成为一件确定的事。

当然也要把话说全：阿司匹林并非没有副作用（比如对胃肠道的刺激、增加出血风险，儿童在某些病毒感染后服用还有罕见但严重的风险），吃不吃、吃多少，要交给医生和说明书，考虑个体差异——这本书不给任何人开处方。科学的诚实是双向的：既不迷信“天然”，也不神化“药片”。

顺便回应一下生活里的同类问题：香草冰淇淋里的香兰素大多是合成的，和香草荚里的主要香气分子是同一种分子；牛仔裤的靛蓝色过去从植物里提取，现在几乎全是合成的，染料分子本身并无不同。你可以基于价格、风味层次、环境代价去偏好天然或合成来源，那都是合理的讨论；但“天然所以安全、合成所以有害”不在合理清单上。

\section{细胞：不用烧瓶的合成大师}

还记得本章埋下的那个痛处吗——化学家很难让反应只发生在想要的位置、只生成想要的那只“手”，只能靠保护基和美纹纸硬扛？

现在告诉你一个让人不太舒服的事实：你身体里的每一个细胞，此刻都在做合成，而且做得比任何化学家都好。在 $37\,^{\circ}\mathrm{C}$、常压、水溶液里，细胞合成蛋白质、脂肪、糖，几十个官能团挤在同一个分子里，反应想动哪个就动哪个，手性从来不出错，不用一滴有机溶剂，不用一个保护基，废物几乎为零。

细胞是怎么做到的？靠一种蛋白质做的“分子机床”——酶。酶把底物分子按精确的角度和位置固定在自己的口袋里，让反应\textbf{只能}在那一个位置、以那一种立体方式发生。化学家梦寐以求的选择性，是酶的日常操作。事实上，制药工业早就“投降”并拜了师：青霉素的半合成模式——让微生物造底盘、化学家做装修——就是把细胞的合成能力和图纸思维嫁接在一起。

但这引出了更深的问题：酶凭什么认位置？细胞这台合成机器的能量从哪里来、又是谁在给它“下图纸”？这些问题超出了有机化学的地界——它们属于下一篇。第六篇，我们从“生命是一种怎样的化学系统”（第 45 章）问起，一路走到糖、脂肪、蛋白质，然后在第 49 章正式拜访酶——你会看到，选择性这个让化学家头秃的难题，生命在几十亿年前就用蛋白质和进化解开了。

\begin{tiaozhan}
想再算深一层吗？（跳过不影响主线。）前面算过：十步线性路线，每步产率 $90\%$，总产率只有 $0.9^{10}\approx 35\%$。但如果把目标分子拆成两个差不多大的片段，\textbf{两支路线各走五步，最后再汇合一步}呢？每支的总产率是 $0.9^5\approx 59\%$，汇合后再乘上最后一步的 $90\%$：$0.9^5\times 0.9\approx 53\%$——同样是十步、同样的单步产率，总产率从 $35\%$ 跳到 $53\%$。秘诀在于：线性路线里，第八步的失败会浪费前七步积累的所有物料；而分支路线里，一支搞砸了只需重做那一支。这叫\textbf{收敛式合成}，是路线设计里仅次于“减少步数”的第二大招。真实药物分子往往被刻意逆合成拆成几个大小相当的片段，正是为了吃到这个数学红利。你可以自己验证：把十步分成“3+3+4”三支再汇合，总产率会更高还是更低？
\end{tiaozhan}

\section*{练一练}

\begin{sikaoti}
\item 第 40 章讲过：羧酸和醇反应生成酯和水。请用逆合成的思路“拆”乙酸乙酯（\chem{CH_3COOCH_2CH_3}，指甲油溶剂和水果香气里都有它）：画出它的酯键位置，标出第一刀切在哪里，写出拆出的两种原料。箭头方向用“实线=合成、虚线=倒推”两种画法各画一遍。
\item 一条八步的制药路线，每步产率 $80\%$。总产率是多少？如果想得到 \ud{100}{g} 最终产品（粗略假设各步分子量相近），第一步大约要投多少原料？算完说说你对“步数”的感受。
\item 原子经济性：第 38 章的加成反应（两个反应物合成一个产物，没有副产物）原子经济性是多少？取代反应（一个原子团换掉另一个原子团）的原子经济性可能达到 $100\%$ 吗？为什么？
\item 画一幅“物质流图”展示保护基的代价：原料 A（带两个官能团）→ 戴保护帽 → 反应 → 摘帽子 → 产物，在每条箭头旁标注可能损失物料的出口。然后回答：为什么说保护基是“必要的浪费”？“必要”和“浪费”各自指什么？
\item 药店的维生素 \chem{C} 片是工厂合成的，橙子里的维生素 \chem{C} 是植物合成的。两者是同一种分子吗？它们在水里的溶解性、被人体吸收的方式会有差别吗？用本章的“分子不记出身”说明理由，再说说“天然维生素 \chem{C} 更好吸收”这类广告词错在哪一层。
\item 你想把一种醇氧化成醛（第 40 章），但反应容易“刹不住车”继续氧化成羧酸。从本章讲的“反应到哪一步”这一层选择性出发，提出至少两条控制思路，并说明每条思路利用了什么原理。
\end{sikaoti}
